%% =====================================================================
%% Chapter 1 -- Introduction (author's rewrite)
%% Every number below comes from the stored runs; see
%% cxr_gui/thesis/PROJECT_DOSSIER.md for its source.
%% =====================================================================
\chapter{Introduction}
\label{ch:introduction}

The automatic generation of reports for chest radiographs has moved from free-text
reporting to grounded reporting in recent years, where findings are described with text
and the location of each finding with a bounding box. This thesis investigates the
possibility of improving upon a grounded model if it goes wrong. Grounded models are
concrete in nature, so when mistakes occur they can be corrected. The question
investigated is whether the corrections can be used to improve the model. To ensure
any new model will not perform worse than the previous, it must be demonstrated that
it will not perform worse than the existing model on unseen data.

\section{Clinical and Technical Background}
\label{sec:perspective}

Chest radiography is the most commonly performed imaging examination
worldwide~\cite{irvin2019chexpert} whilst the simple chest radiograph is the most
frequently requested examination for pulmonary disorders~\cite{raoof2012chest}. Every
examination must be read and reported. Models to help generate reports have been
developed. However, it is difficult to verify whether free texts are accurate as the
reader cannot see where each sentence refers to on the image. There will be no
difference between hallucinated findings and correct findings.

This issue can be overcome by using the grounded report generation approach to
generate reports. Within this field the ``grounded report generation'' task has been
developed. This involves generating a chest X-ray report with the localisation of
individual findings, and was introduced with MAIRA-2~\cite{bannur2024maira2}. Another
example of this is CURE~\cite{messina2026cure}. This uses a low rank adapter
(LoRA~\cite{hu2021lora}) on top of a publicly available vision-language medical model
(MedGemma~\cite{sellergren2025medgemma}). This model learns to localise and describe
anatomy, including normal anatomy, to reduce hallucinated abnormalities. Both models
were trained on the grounded bilingual dataset PadChest-GR~\cite{castro2024padchestgr}.

However, the models are very different from a technical perspective. MAIRA-2 is a large
model whereas CURE is a small number of adapter weights on a public model. This means
that CURE can be further trained on just one workstation class GPU. This means it can
be version controlled. This is what allows a correction loop. However, this can be
risky as the model may forget what it learned before and suffer from catastrophic
forgetting~\cite{kirkpatrick2017forgetting}. It may perform worse on data outside of
the training data. Therefore it is important to have a quantifiable and repeatable
method of verifying any new models can be safely updated.

\section{Problem Statement}
\label{sec:problem}

There are two problems presented.

Firstly, evaluation. If a model is evaluated on data it has been trained on this
measures how well it remembers the data rather than how well it generalises. This is a
common cause of exaggerated results in machine learning based
research~\cite{kapoor2022leakage}. In this case, both MAIRA-2 and CURE were trained on
PadChest-GR. Therefore any comparison must take into account the official split. The
first benchmark created did not do this. 129 out of 200 images were from the training
split. This made CURE appear to outperform MAIRA-2. However, when limiting the data to
the held out test images the ranking was reversed (Section~\ref{sec:leakage}).
Therefore if one grounded model performs better than another either internally or
published, it cannot be trusted until it has been measured on studies that neither
model has seen.

Secondly, improvements. When reviewing a grounded report, a reviewer provides the exact
information needed to improve a model. They know what boxes were correct and incorrect
and what findings were hallucinated and missed. However there must be a quantifiable
and repeatable way to update the model. This should take the corrections and use them
as training data in the same format as the model's output. Also, this method should
protect the model's knowledge and decide whether the new model can replace the old
based on held out data using a statistical test.

\section{Research Objectives}
\label{sec:objectives}

The following four research questions will be investigated.

\begin{description}[leftmargin=0pt, itemsep=4pt]
  \item[RQ1.] What is the effect of evaluating CURE and MAIRA-2 on the held out test
        studies from PadChest-GR? How do they perform on the grounded detection task
        (F1 score with IoU~$\ge$~0.5, hallucinated findings) and finding name task
        (keyword F1) whilst respecting the official split?
  \item[RQ2.] Can CURE be improved by continuing to train the CURE adapter on
        corrections, when each new version must pass an adoption gate on the held out
        test set? Does this affect the relationship with MAIRA-2?
  \item[RQ3.] What effect does the number of corrections have? (50, 150 and 300).
  \item[RQ4 (Exploratory).] Does the text generated by converting corrected findings to
        a narrative by a general purpose language model stay within the findings
        provided? This is only an automated consistency check.
\end{description}

In order to answer RQ2 and RQ3 an oracle reviewer was used to generate corrections.
This automatically generates corrections from the ground truth of the PadChest-GR
validation split (Section~\ref{sec:oracle}). This reviewer can see all errors and never
makes an error. This means the results are an upper bound of what can be achieved.

All experiments were conducted on one dataset, PadChest-GR. The official split was
followed with the validation split providing the corrections, the training split
providing replay examples and the test split being used for evaluation. The findings
detection and naming task was evaluated on the ground truth of the dataset, however no
clinical evaluation was carried out. Only an automated consistency check was conducted
on the narrative reports (RQ4). The reviewer application was designed to allow real
corrections to be used but all corrections come from the oracle.

\section{Contributions}
\label{sec:contributions}

The main contributions of this thesis are:

\begin{enumerate}[label=(\arabic*), leftmargin=*, itemsep=4pt]
  \item \textbf{A leakage controlled benchmark comparing CURE and MAIRA-2.} Both
        models were scored using one scoring system on all 604 test studies with ground
        truth boxes on PadChest-GR. CURE performs significantly worse than MAIRA-2 on
        detection F1 on the held out data (Section~\ref{sec:baseline}).
  \item \textbf{A reviewer application to capture the corrections.} A desktop
        application was created to allow the reviewer to move, rename, add and remove
        boxes. All corrections including removals were exported in a format ready to use
        for training. A server was used to run the model, store corrections and manage
        adapter versions (Sections~\ref{sec:app} and~\ref{sec:server}).
  \item \textbf{A gated continual learning loop.} New adapter versions were trained
        using corrections and replay examples from the training split. Each new version
        must pass an adoption gate using regression tolerances and a paired bootstrap
        test on the held out test set before it can be used
        (Sections~\ref{sec:training} and~\ref{sec:gate}).
  \item \textbf{Learning curve.} All versions trained with 50, 150 and 300 oracle
        corrections passed the adoption gate. Every version had a higher precision and a
        lower rate of hallucinated findings than the original model. After 150
        corrections the difference in detection F1 per image between CURE and MAIRA-2
        was no longer statistically significant (Sections~\ref{sec:curve}
        and~\ref{sec:significance}).
  \item \textbf{Exploratory narrative check.} A second model was used to check whether
        the reports generated by the language model contain any unsupported statements.
        This is an automated non-clinical check and the limitations are discussed
        (Section~\ref{sec:report-check}).
\end{enumerate}

\section{Structure}
\label{sec:structure}

The remainder of this thesis is structured as follows. Chapter~\ref{ch:background}
provides an overview of the background information. This includes how image models
have evolved to vision-language models, how a vision-language model is built, adapting
a base model, the grounded report generation models used for comparison, the
PadChest-GR dataset and the evaluation methodology. Chapter~\ref{ch:methods} covers
the materials and methods of the research. This includes the dataset and splits,
reviewer application, server, oracle reviewer, continual learning and adoption gate.
Chapter~\ref{ch:results} provides the experiments carried out. This includes the
baseline comparison between CURE and MAIRA-2, learning curve, statistical analysis and
exploratory narrative check. Chapter~\ref{ch:conclusion} is a summary of the findings,
discussion of limitations and future directions. Appendix~\ref{app:provenance} shows a
mapping of all results reported to the output they come from.
Appendix~\ref{app:reproduction} provides details of the software environment and
commands used to recreate the experiments.
